\section{Finite sets}

\import{set/cons.tex}
\import{set/powerset.tex}
\import{function.tex}
\import{nat.tex}

The finite sets are the sets enumerated without repetition by a von Neumann
natural number.  In particular, this definition is independent of any later
construction of integers or real numbers.

\begin{proposition}\label{powerset_cons_image}
    Suppose $g$ is a function on $\pow{A}$.
    Suppose for all $B\in\pow{A}$ we have $g(B)=\cons{b}{B}$.
    Then $\pow{\cons{b}{A}}=\pow{A}\union\img{g}{\pow{A}}$.
\end{proposition}
\begin{proof}
    Show that $\pow{\cons{b}{A}}\subseteq
        \pow{A}\union\img{g}{\pow{A}}$.
    \begin{subproof}
        It suffices to show that every element of $\pow{\cons{b}{A}}$ is
            an element of $\pow{A}\union\img{g}{\pow{A}}$
            by \obv{subseteq}.
        Fix $C\in\pow{\cons{b}{A}}$.
        $C\subseteq\cons{b}{A}$ by \obv{pow_iff,subseteq}.
        \begin{byCase}
        \caseOf{$b\in C$.}
            Let $B=C\setminus\{b\}$.
            $B\subseteq A$ by \auto{subseteq,setminus,cons_iff}.
            Thus $B\in\pow{A}$ by \obv{pow_iff,subseteq}.
            Hence $B\in\dom{g}$.
            $C=\cons{b}{B}$ by \auto{cons_remove}.
            Hence $C=g(B)$.
            Thus $C\in\img{g}{\pow{A}}$
                by \auto{img_of_function_intro,inter_intro}.
            Follows by \auto{union_iff}.
        \caseOf{$b\notin C$.}
            $C\subseteq A$ by \auto{subseteq,cons_iff}.
            Thus $C\in\pow{A}$ by \obv{pow_iff,subseteq}.
            Follows by \auto{union_iff}.
        \end{byCase}
    \end{subproof}

    Show that $\pow{A}\union\img{g}{\pow{A}}
        \subseteq\pow{\cons{b}{A}}$.
    \begin{subproof}
        It suffices to show that every element of
            $\pow{A}\union\img{g}{\pow{A}}$ is an element of
            $\pow{\cons{b}{A}}$ by \obv{subseteq}.
        Fix $C\in\pow{A}\union\img{g}{\pow{A}}$.
        \begin{byCase}
        \caseOf{$C\in\pow{A}$.}
            $C\subseteq A$ by \obv{pow_iff,subseteq}.
            Thus $C\subseteq\cons{b}{A}$ by \auto{subseteq,cons_iff}.
            Follows by \obv{pow_iff,subseteq}.
        \caseOf{$C\notin\pow{A}$.}
            $C\in\img{g}{\pow{A}}$ by \auto{union_iff}.
            Take $B\in\pow{A}$ such that $C=g(B)$
                by \auto{img_of_function_elim,inter_elim_right}.
            $C=\cons{b}{B}$.
            $B\subseteq A$ by \obv{pow_iff,subseteq}.
            Thus $C\subseteq\cons{b}{A}$
                by \auto{subseteq,cons_iff}.
            Follows by \obv{pow_iff,subseteq}.
        \end{byCase}
    \end{subproof}
    Follows by \obv{subseteq_antisymmetric}.
\end{proof}

\begin{definition}\label{finite}
    $X$ is finite iff there exists a natural number $n$ such that
        there exists a bijection from $n$ to $X$.
\end{definition}

\begin{abbreviation}\label{infinite}
    $X$ is infinite iff $X$ is not finite.
\end{abbreviation}

\begin{definition}\label{finite_subsets}
    $\finiteSubsets{X}=\{F\in\pow{X}\mid \text{$F$ is finite}\}$.
\end{definition}

\begin{definition}\label{infinite_subsets}
    $\infiniteSubsets{X}=\{A\in\pow{X}\mid \text{$A$ is infinite}\}$.
\end{definition}

\begin{proposition}\label{finite_subsets_iff}
    $F\in\finiteSubsets{X}$ iff $F\subseteq X$ and $F$ is finite.
\end{proposition}
\begin{proof}
    Follows by \auto{finite_subsets,pow_iff,subseteq}.
\end{proof}

\begin{proposition}\label{infinite_subsets_iff}
    $A\in\infiniteSubsets{X}$ iff $A\subseteq X$ and $A$ is infinite.
\end{proposition}
\begin{proof}
    Follows by \auto{infinite_subsets,pow_iff,subseteq}.
\end{proof}

\begin{proposition}\label{finite_subsets_subset_powerset}
    $\finiteSubsets{X}\subseteq\pow{X}$.
\end{proposition}
\begin{proof}
    Follows by \obv{finite_subsets_iff,pow_iff,subseteq}.
\end{proof}

\begin{proposition}\label{infinite_subsets_subset_powerset}
    $\infiniteSubsets{X}\subseteq\pow{X}$.
\end{proposition}
\begin{proof}
    Follows by \obv{infinite_subsets_iff,pow_iff,subseteq}.
\end{proof}

\begin{classification}\label{emptyset_is_finite}
    $\emptyset$ is finite.
\end{classification}
\begin{proof}
    $\identity{\emptyset}\in\bijections{\emptyset}{\emptyset}$
        by \obv{id_elem_bijections}.
    Follows by \obv{finite,omega_empty}.
\end{proof}

\begin{proposition}\label{infinite_is_inhabited}
    Suppose $X$ is infinite.
    Then $X$ is inhabited.
\end{proposition}
\begin{proof}
    Suppose not.
    Hence $X=\emptyset$ by \obv{nonempty_inhabited}.
    $\emptyset$ is finite by \obv{emptyset_is_finite}.
    Hence $X$ is finite.
    Contradiction.
\end{proof}

\begin{proposition}\label{bijection_cons}
    Suppose $f\in\bijections{A}{B}$.
    Suppose $a\notin A$ and $b\notin B$.
    Then $\cons{(a,b)}{f}\in\bijections{\cons{a}{A}}{\cons{b}{B}}$.
\end{proposition}
\begin{proof}
    Follows by \auto{bijections_dom,dom_intro,bijection_is_function,relation,cons_iff,rightunique,pair_eq_iff,function_rightunique,ran_of_surj,dom_cons,ran_cons,function_apply_elim,ran_intro,funs_intro,funs_surj_iff,injective,bijections}.
\end{proof}

\begin{proposition}\label{finite_cons}
    Suppose $X$ is finite.
    Then $\cons{x}{X}$ is finite.
\end{proposition}
\begin{proof}
    Take $n$ such that $n$ is a natural number and
        there exists $f\in\bijections{n}{X}$ by \auto{finite}.
    Take $f$ such that $f\in\bijections{n}{X}$ by \auto.
    \begin{byCase}
    \caseOf{$x\in X$.}
        $\cons{x}{X}=X$ by \obv{cons_absorb}.
        Thus $\cons{x}{X}$ is finite.
    \caseOf{$x\notin X$.}
        $n\notin n$ by \obv{in_irrefl}.
        $\cons{(n,x)}{f}\in\bijections{\cons{n}{n}}{\cons{x}{X}}$
            by \obv{bijection_cons}.
        $\cons{n}{n}=\suc{n}$ by \obv{suc,cons_iff,setext}.
        $\suc{n}$ is a natural number by \obv{omega_successor}.
        Thus $\cons{x}{X}$ is finite by \auto{finite}.
    \end{byCase}
\end{proof}

\begin{corollary}\label{singleton_is_finite}
    $\{x\}$ is finite.
\end{corollary}
\begin{proof}
    $\cons{x}{\emptyset}$ is finite
        by \obv{emptyset_is_finite,finite_cons}.
    Every element of $\cons{x}{\emptyset}$ is an element of $\{x\}$
        by \obv{cons_iff,singleton_iff}.
    Every element of $\{x\}$ is an element of $\cons{x}{\emptyset}$
        by \obv{cons_iff,singleton_iff}.
    Hence $\cons{x}{\emptyset}=\{x\}$ by set extensionality.
\end{proof}

\begin{corollary}\label{pair_is_finite}
    $\{x,y\}$ is finite.
\end{corollary}
\begin{proof}
    $\cons{y}{\emptyset}$ is finite
        by \obv{emptyset_is_finite,finite_cons}.
    Thus $\cons{x}{\cons{y}{\emptyset}}$ is finite by \obv{finite_cons}.
    $\{x,y\}=\cons{x}{\cons{y}{\emptyset}}$
        by \auto{setext,pairset_iff,cons_iff,emptyset}.
\end{proof}

\begin{proposition}\label{bijection_remove_successor}
    Suppose $f\in\bijections{\suc{n}}{Y}$.
    Then $\restrl{f}{n}\in\bijections{n}{Y\setminus\{f(n)\}}$.
\end{proposition}
\begin{proof}
    Let $g=\restrl{f}{n}$.
    $f$ is a function by \obv{bijection_is_function}.
    $f$ is injective by \obv{bijections}.
    $\dom{f}=\suc{n}$ by \obv{bijections_dom}.
    $\ran{f}=Y$ by \obv{bijections,ran_of_surj}.
    $n\subseteq\suc{n}$ by \obv{subseteq_self_suc_intro}.
    $g$ is a function by \obv{restrl_of_function_is_function}.
    $\dom{g}=n$ by \auto{dom_of_restrl_eq_inter,inter_absorb_supseteq_left}.

    Show that $\ran{g}\subseteq Y\setminus\{f(n)\}$.
    \begin{subproof}
        It suffices to show that every element of $\ran{g}$ is an element
            of $Y\setminus\{f(n)\}$ by \obv{subseteq}.
        Fix $z\in\ran{g}$.
        Take $a$ such that $(a,z)\in g$ by \obv{ran_iff}.
        Then $(a,z)\in f$ and $a\in n$ by \obv{restrl_iff}.
        Thus $z\in Y$ by \obv{ran_intro}.
        Show that $z\neq f(n)$.
        \begin{subproof}
            Suppose not.
            $n\in\dom{f}$ by \obv{suc_intro_self}.
            Thus $(n,f(n))\in f$ by \obv{function_apply_elim}.
            Hence $a=n$ by \auto{injective}.
            Contradiction by \obv{in_irrefl}.
        \end{subproof}
        Follows by \auto{setminus_intro,singleton_iff}.
    \end{subproof}

    Show that $Y\setminus\{f(n)\}\subseteq\ran{g}$.
    \begin{subproof}
        It suffices to show that every element of $Y\setminus\{f(n)\}$ is
            an element of $\ran{g}$ by \obv{subseteq}.
        Fix $z\in Y\setminus\{f(n)\}$.
        Then $z\neq f(n)$ by \obv{setminus_elim_right,singleton_iff}.
        $z\in Y$ by \obv{setminus_elim_left}.
        Take $a\in\dom{f}$ such that $f(a)=z$ by \auto{fun_ran_iff}.
        $a=n$ or $a\in n$ by \obv{suc_iff}.
        Thus $a\in n$.
        $(a,z)\in f$ by \auto{function_apply_elim}.
        Hence $(a,z)\in g$ by \obv{restrl_iff}.
        Follows by \obv{ran_intro}.
    \end{subproof}
    Hence $\ran{g}=Y\setminus\{f(n)\}$ by \obv{subseteq_antisymmetric}.
    For all $a\in\dom{g}$ we have $g(a)\in Y\setminus\{f(n)\}$
        by \auto{function_apply_elim,ran_intro}.
    Thus $g$ is a function to $Y\setminus\{f(n)\}$.
    Hence $g\in\funs{n}{Y\setminus\{f(n)\}}$ by \obv{funs_intro}.
    Thus $g\in\surj{n}{Y\setminus\{f(n)\}}$ by \obv{funs_surj_iff}.
    $g$ is injective by \obv{restrl_injective}.
    Follows by \obv{bijections}.
\end{proof}

\begin{theorem}\label{finite_induction}
    Suppose $X$ is finite.
    Suppose $\emptyset\in C$.
    Suppose for all $Y,y$ such that $Y\in C$ and $y\in X$ we have
        $\cons{y}{Y}\in C$.
    Then $X\in C$.
\end{theorem}
\begin{proof}
    Let $S=\{n\in\naturals\mid
        \text{for all $Y,f$ such that $Y\subseteq X$ and
            $f\in\bijections{n}{Y}$ we have $Y\in C$}\}$.

    Show that $\emptyset\in S$.
    \begin{subproof}
        $\emptyset\in\naturals$ by \obv{omega_empty}.
        For all $Y,f$ such that $Y\subseteq X$ and
            $f\in\bijections{\emptyset}{Y}$ we have $Y=\emptyset$
            by \auto{bijections,surj,emptyset,setext}.
    \end{subproof}

    Show that for all $n\in S$ we have $\suc{n}\in S$.
    \begin{subproof}
        Fix $n\in S$.
        Show that for all $Y,f$ such that $Y\subseteq X$ and
            $f\in\bijections{\suc{n}}{Y}$ we have $Y\in C$.
        \begin{subproof}
            Fix $Y$.
            Fix $f$.
            Assume $Y\subseteq X$ and $f\in\bijections{\suc{n}}{Y}$.
            $f$ is a function by \obv{bijection_is_function}.
            $\dom{f}=\suc{n}$ by \obv{bijections_dom}.
            Let $y=f(n)$.
            $n\in\dom{f}$ by \obv{suc_intro_self}.
            Thus $y\in Y$ by \auto{function_apply_elim,ran_intro,bijections,ran_of_surj}.
            Hence $y\in X$ by \obv{elem_subseteq}.
            $\restrl{f}{n}\in\bijections{n}{Y\setminus\{y\}}$
                by \auto{bijection_remove_successor}.
            $Y\setminus\{y\}\subseteq X$
                by \obv{setminus_subseteq,subseteq_transitive}.
            Thus $Y\setminus\{y\}\in C$.
            Hence $\cons{y}{Y\setminus\{y\}}\in C$.
            $\cons{y}{Y\setminus\{y\}}=Y$ by \obv{cons_remove}.
            Thus $Y\in C$.
        \end{subproof}
        Follows by \obv{omega_successor}.
    \end{subproof}

    Thus $S$ is an inductive set by \obv{inductive_set}.
    Take $n$ such that $n$ is a natural number and
        there exists $f\in\bijections{n}{X}$ by \auto{finite}.
    Take $f$ such that $f\in\bijections{n}{X}$ by \auto.
    $n\in S$ by \obv{omega_minimal}.
    Follows by \auto{subseteq_refl}.
\end{proof}

\begin{proposition}\label{subseteq_finite_is_finite}
    Suppose $A\subseteq B$.
    Suppose $B$ is finite.
    Then $A$ is finite.
\end{proposition}
\begin{proof}
    Let $C=\{Y\in\pow{B}\mid
        \text{for all $D$ such that $D\subseteq Y$ we have $D$ is finite}\}$.
    Thus $A$ is finite by \auto{powerset_bottom,subseteq_emptyset_iff,emptyset_is_finite,powerset_elim,cons_subseteq_intro,powerset_intro,subseteq_cons_elim,finite_cons,cons_remove,finite_induction}.
\end{proof}

\begin{theorem}\label{infinite_superset}
    Suppose $A\subseteq B$ and $A$ is infinite.
    Then $B$ is infinite.
\end{theorem}
\begin{proof}
    Suppose not.
    Then $A$ is finite by \obv{subseteq_finite_is_finite}.
    Contradiction.
\end{proof}

\begin{corollary}\label{finite_intersection}
    Suppose $A$ is finite or $B$ is finite.
    Then $A\inter B$ is finite.
\end{corollary}
\begin{proof}
    Follows by \auto{inter_lower_left,inter_lower_right,subseteq_finite_is_finite}.
\end{proof}

\begin{proposition}\label{union_of_finite_is_finite}
    Suppose $A$ is finite and $B$ is finite.
    Then $A\union B$ is finite.
\end{proposition}
\begin{proof}
    Let $C=\{Y\in\pow{A}\mid \text{$Y\union B$ is finite}\}$.
    Thus $A\union B$ is finite by \auto{powerset_bottom,union_comm,union_emptyset,powerset_elim,cons_subseteq_intro,powerset_intro,finite_cons,union_cons,finite_induction}.
\end{proof}

\begin{theorem}\label{infinite_minus_finite_is_infinite}
    Suppose $A$ is infinite and $F$ is finite.
    Then $A\setminus F$ is infinite.
\end{theorem}
\begin{proof}
    Suppose not.
    $A\inter F$ is finite by \obv{finite_intersection}.
    $A=(A\setminus F)\union(A\inter F)$
        by \auto{setext,setminus,inter,union_iff}.
    Hence $A$ is finite by \obv{union_of_finite_is_finite}.
    Contradiction.
\end{proof}

\begin{corollary}\label{infinite_minus_finite_inhabited}
    Suppose $A$ is infinite and $F$ is finite.
    Then $A\setminus F$ is inhabited.
\end{corollary}
\begin{proof}
    $A\setminus F$ is infinite by \obv{infinite_minus_finite_is_infinite}.
    Suppose not.
    Then $A\setminus F=\emptyset$ by \obv{nonempty_inhabited}.
    Hence $A\setminus F$ is finite by \obv{emptyset_is_finite}.
    Contradiction.
\end{proof}

\begin{theorem}\label{finite_union_has_finite_trace}
    Suppose $F$ is finite.
    Suppose for every $U\in F$ we have $U\inter D$ is finite.
    Then $\unions{F}\inter D$ is finite.
\end{theorem}
\begin{proof}
    Let $C=\{G\in\pow{F}\mid
        \text{$\unions{G}\inter D$ is finite}\}$.
    $\emptyset\subseteq F$ by \obv{emptyset_subseteq}.
    Hence $\emptyset\in\pow{F}$ by \obv{pow_iff,subseteq}.
    \begin{align*}
        \unions{\emptyset}\inter D
            &= \emptyset\inter D
                \explanation{by \obv{unions_emptyset}}
            \\
            &= D\inter\emptyset
                \explanation{by \obv{inter_comm}}
            \\
            &= \emptyset
                \explanation{by \obv{inter_emptyset}}
    \end{align*}
    Hence $\unions{\emptyset}\inter D$ is finite
        by \obv{emptyset_is_finite}.
    Thus $\emptyset\in C$.
    Show that for all $G,U$ such that $G\in C$ and $U\in F$ we have
        $\cons{U}{G}\in C$.
    \begin{subproof}
        Fix $G,U$.
        Assume $G\in C$ and $U\in F$.
        $G\subseteq F$ by \obv{pow_iff,subseteq}.
        Hence $\cons{U}{G}\subseteq F$ by \obv{cons_subseteq_intro}.
        Hence $\cons{U}{G}\in\pow{F}$ by \obv{pow_iff,subseteq}.
        $U\inter D$ is finite.
        $\unions{G}\inter D$ is finite.
        \begin{align*}
            \unions{\cons{U}{G}}\inter D
                &= (U\union\unions{G})\inter D
                    \explanation{by \obv{unions_cons}}
                \\
                &= (U\inter D)\union(\unions{G}\inter D)
                    \explanation{by \auto{inter_distrib_union,inter_comm}}
        \end{align*}
        Hence $\unions{\cons{U}{G}}\inter D$ is finite
            by \obv{union_of_finite_is_finite}.
    \end{subproof}
    $F\in C$ by \obv{finite_induction,subseteq_refl,pow_iff}.
\end{proof}

\begin{proposition}\label{unions_of_finite_family_is_finite}
    Suppose $F$ is finite.
    Suppose every element of $F$ is finite.
    Then $\unions{F}$ is finite.
\end{proposition}
\begin{proof}
    Let $C=\{G\in\pow{F}\mid \text{$\unions{G}$ is finite}\}$.
    $\emptyset\in C$ by \obv{powerset_bottom,unions_emptyset,emptyset_is_finite}.
    For all $G,X$ such that $G\in C$ and $X\in F$ we have
        $\cons{X}{G}\in C$
        by \auto{powerset_elim,cons_subseteq_intro,powerset_intro,unions_cons,union_of_finite_is_finite}.
    Follows by \auto{finite_induction}.
\end{proof}

\begin{theorem}\label{finite_family_union_in_binary_union_closed_family}
    Suppose $\emptyset\in F$.
    Suppose for all $A,C\in F$ we have $A\union C\in F$.
    Suppose $H\subseteq F$ and $H$ is finite.
    Then $\unions{H}\in F$.
\end{theorem}
\begin{proof}
    Let $C=\{G\in\pow{F}\mid \unions{G}\in F\}$.
    $\emptyset\in\pow{F}$ by \obv{powerset_bottom}.
    $\unions{\emptyset}=\emptyset$ by \obv{unions_emptyset}.
    Hence $\emptyset\in C$.
    Show that for all $G,A$ such that $G\in C$ and $A\in H$ we have
        $\cons{A}{G}\in C$.
    \begin{subproof}
        Fix $G,A$.
        Assume $G\in C$ and $A\in H$.
        $G\subseteq F$ by \obv{pow_iff,subseteq}.
        $A\in F$ by \obv{subseteq}.
        Hence $\cons{A}{G}\subseteq F$ by
            \obv{cons_subseteq_intro}.
        Thus $\cons{A}{G}\in\pow{F}$ by
            \obv{powerset_intro}.
        $\unions{G}\in F$.
        Hence $A\union\unions{G}\in F$.
        $\unions{\cons{A}{G}}=A\union\unions{G}$ by
            \obv{unions_cons}.
    \end{subproof}
    Hence $H\in C$ by \obv{finite_induction}.
\end{proof}

\begin{proposition}\label{finite_function_image}
    Suppose $f$ is a function.
    Suppose $X$ is finite and $X\subseteq\dom{f}$.
    Then $\img{f}{X}$ is finite.
\end{proposition}
\begin{proof}
    Let $C=\{Y\in\pow{X}\mid \text{$\img{f}{Y}$ is finite}\}$.
    $\emptyset\in C$ by \obv{powerset_bottom,img_emptyset,emptyset_is_finite}.
    Show that for all $Y,x$ such that $Y\in C$ and $x\in X$ we have
        $\cons{x}{Y}\in C$.
    \begin{subproof}
        Fix $Y$.
        Fix $x$.
        Assume $Y\in C$ and $x\in X$.
        $x\in\dom{f}$ by \obv{elem_subseteq}.
        $\img{f}{\{x\}}=\{f(x)\}$
            by \auto{setext,img_iff,singleton_iff,function_apply_iff}.
        $\img{f}{\cons{x}{Y}}=\cons{f(x)}{\img{f}{Y}}$
            by \auto{union_eq_cons,img_union}.
        $\img{f}{Y}$ is finite.
        Thus $\img{f}{\cons{x}{Y}}$ is finite by \auto{finite_cons}.
        $\cons{x}{Y}\in\pow{X}$
            by \obv{powerset_elim,cons_subseteq_intro,powerset_intro}.
    \end{subproof}
    Follows by \auto{finite_induction}.
\end{proof}

\begin{classification}\label{finite_powerset}
    Suppose $X$ is finite.
    Then $\pow{X}$ is finite.
\end{classification}
\begin{proof}
    Let $C=\{Y\in\pow{X}\mid \text{$\pow{Y}$ is finite}\}$.
    $\pow{\emptyset}$ is finite
        by \auto{powerset_emptyset,singleton_is_finite}.
    Thus $\emptyset\in C$ by \obv{powerset_bottom}.

    Show that for all $Y,y$ such that $Y\in C$ and $y\in X$ we have
        $\cons{y}{Y}\in C$.
    \begin{subproof}
        Fix $Y$.
        Fix $y$.
        Assume $Y\in C$ and $y\in X$.
        $Y\subseteq X$ by \obv{pow_iff,subseteq}.
        $\pow{Y}$ is finite.

        Let $g=\{(B,\cons{y}{B})\mid B\in\pow{Y}\}$.
        $g$ is a relation by \auto{relation}.
        $g$ is right-unique by \auto{rightunique,pair_eq_iff}.
        $\dom{g}=\pow{Y}$ by \auto{setext,dom_iff,pair_eq_iff}.
        Hence $g$ is a function on $\pow{Y}$.
        For all $B\in\pow{Y}$ we have $g(B)=\cons{y}{B}$
            by \obv{function_apply_intro}.

        $\img{g}{\pow{Y}}$ is finite
            by \obv{finite_function_image,subseteq_refl}.
        $\pow{\cons{y}{Y}}=\pow{Y}\union\img{g}{\pow{Y}}$
            by \obv{powerset_cons_image}.
        Thus $\pow{\cons{y}{Y}}$ is finite
            by \obv{union_of_finite_is_finite}.
        $\cons{y}{Y}\subseteq X$ by \obv{cons_subseteq_intro}.
        Hence $\cons{y}{Y}\in\pow{X}$ by \obv{pow_iff,subseteq}.
    \end{subproof}
    $X\in C$ by \obv{finite_induction}.
\end{proof}

\begin{proposition}\label{times_cons_left}
    $\cons{x}{X}\times Y=(\{x\}\times Y)\union(X\times Y)$.
\end{proposition}
\begin{proof}
    Every element of $\cons{x}{X}\times Y$ is an element of
        $(\{x\}\times Y)\union(X\times Y)$
        by \auto{times,cons_iff,singleton_iff,union_iff,pair_eq_iff}.
    Every element of $(\{x\}\times Y)\union(X\times Y)$ is an element of
        $\cons{x}{X}\times Y$
        by \auto{times,cons_iff,singleton_iff,union_iff,pair_eq_iff}.
    Follows by set extensionality.
\end{proof}

\begin{proposition}\label{times_cons_right}
    $X\times\cons{y}{Y}=(X\times\{y\})\union(X\times Y)$.
\end{proposition}
\begin{proof}
    Every element of $X\times\cons{y}{Y}$ is an element of
        $(X\times\{y\})\union(X\times Y)$
        by \auto{times,cons_iff,singleton_iff,union_iff,pair_eq_iff}.
    Every element of $(X\times\{y\})\union(X\times Y)$ is an element of
        $X\times\cons{y}{Y}$
        by \auto{times,cons_iff,singleton_iff,union_iff,pair_eq_iff}.
    Follows by set extensionality.
\end{proof}

\begin{proposition}\label{times_singletons}
    $\{x\}\times\{y\}=\{(x,y)\}$.
\end{proposition}
\begin{proof}
    Follows by
        \auto{setext,times,singleton_iff,pair_eq_iff}.
\end{proof}

\begin{lemma}\label{singleton_times_finite}
    Suppose $Y$ is finite.
    Then $\{x\}\times Y$ is finite.
\end{lemma}
\begin{proof}
    Let $C=\{Z\in\pow{Y}\mid
        \text{$\{x\}\times Z$ is finite}\}$.
    $\emptyset\in C$
        by
        \obv{powerset_bottom,times_empty_right,emptyset_is_finite}.
    Show that for all $Z,z$ such that $Z\in C$ and $z\in Y$ we have
        $\cons{z}{Z}\in C$.
    \begin{subproof}
        Fix $Z,z$.
        Assume $Z\in C$ and $z\in Y$.
        $Z\subseteq Y$ by \obv{powerset_elim}.
        Hence $\cons{z}{Z}\subseteq Y$
            by \obv{cons_subseteq_intro}.
        Hence $\cons{z}{Z}\in\pow{Y}$ by \obv{powerset_intro}.
        $\{x\}\times\cons{z}{Z}
            =\cons{(x,z)}{\{x\}\times Z}$
            by
            \auto{times_cons_right,times_singletons,union_eq_cons}.
        $\{x\}\times Z$ is finite.
        Hence $\{x\}\times\cons{z}{Z}$ is finite
            by \auto{finite_cons}.
    \end{subproof}
    $Y\in C$ by \obv{finite_induction}.
\end{proof}

\begin{theorem}\label{finite_times}
    Suppose $X$ is finite and $Y$ is finite.
    Then $X\times Y$ is finite.
\end{theorem}
\begin{proof}
    Let $C=\{Z\in\pow{X}\mid \text{$Z\times Y$ is finite}\}$.
    $\emptyset\in C$
        by
        \obv{powerset_bottom,times_empty_left,emptyset_is_finite}.
    Show that for all $Z,z$ such that $Z\in C$ and $z\in X$ we have
        $\cons{z}{Z}\in C$.
    \begin{subproof}
        Fix $Z,z$.
        Assume $Z\in C$ and $z\in X$.
        $Z\subseteq X$ by \obv{powerset_elim}.
        Hence $\cons{z}{Z}\subseteq X$
            by \obv{cons_subseteq_intro}.
        Hence $\cons{z}{Z}\in\pow{X}$ by \obv{powerset_intro}.
        $\{z\}\times Y$ is finite
            by \obv{singleton_times_finite}.
        $Z\times Y$ is finite.
        $\cons{z}{Z}\times Y
            =(\{z\}\times Y)\union(Z\times Y)$
            by \obv{times_cons_left}.
        Hence $\cons{z}{Z}\times Y$ is finite
            by \obv{union_of_finite_is_finite}.
    \end{subproof}
    $X\in C$ by \obv{finite_induction}.
\end{proof}

\begin{definition}\label{enumeration}
    $f$ is an enumeration of $X$ iff there exists a natural number $n$
        such that $f\in\bijections{n}{X}$.
\end{definition}

\begin{proposition}\label{finite_iff_has_enumeration}
    $X$ is finite iff there exists $f$ such that $f$ is an enumeration of $X$.
\end{proposition}
\begin{proof}
    Follows by \auto{finite,enumeration}.
\end{proof}

\begin{definition}\label{finite_sequence}
    $s$ is a finite sequence in $A$ iff there exists a natural number $n$
        such that $s\in\funs{n}{A}$.
\end{definition}

\begin{proposition}\label{emptyset_finite_sequence}
    $\emptyset$ is a finite sequence in $A$.
\end{proposition}
\begin{proof}
    $\emptyset\in\funs{\emptyset}{A}$
        by \obv{codom_of_emptyset_can_be_anything,dom_emptyset,funs_intro}.
    Follows by \obv{finite_sequence,omega_empty}.
\end{proof}

\begin{proposition}\label{finite_sequence_domain}
    Suppose $s$ is a finite sequence in $A$.
    Then $\dom{s}$ is a natural number and $s\in\funs{\dom{s}}{A}$.
\end{proposition}
\begin{proof}
    Follows by \auto{finite_sequence,funs_elim}.
\end{proof}
